\documentclass[10pt,a4paper]{amsart}
\usepackage[margin=19mm]{geometry}
\usepackage{amsmath,amssymb,mathtools}
\usepackage[scr=rsfs,cal=euler]{mathalfa}
\usepackage{xfrac}
\usepackage[unicode,hidelinks,bookmarks=false]{hyperref}
\newcommand{\ie}{i.e.\ }
\newcommand{\cf}{cf.\ }
\newcommand{\resp}{respectively\ }
\newcommand{\Subsection}{\subsection}
\providecommand{\nobreakdash}{\nobreak\hbox{-}}
\setlength{\emergencystretch}{2em}
\multlinegap0pt
\usepackage{bm}
\usepackage{bbm}
\usepackage{dsfont}
\usepackage{xspace}
\usepackage{cancel}
\usepackage{mathtools}
\usepackage{amsthm}
\makeatletter
\@ifundefined{theorem}{\newtheorem{theorem}{Theorem}}{}
\@ifundefined{lemma}{\newtheorem{lemma}[theorem]{Lemma}}{}
\@ifundefined{observation}{\newtheorem{observation}{Observation}}{}
\makeatother
\title[Openly disjoint cycles and directed tree-width of regular di]{Openly disjoint cycles and directed tree-width of regular digraphs}
\author{Raphael Steiner}
\date{Source: arXiv:2604.13700v2 (CC BY 4.0)}
\begin{document}
\maketitle

\noindent\emph{Source and licence.} Openly disjoint cycles and directed tree-width of regular digraphs. Version arXiv:2604.13700v2 (CC BY 4.0), distributed under the Creative Commons Attribution 4.0 International licence, as recorded in the arXiv licence metadata for this exact version and shown on the abstract page https://arxiv.org/abs/2604.13700v2. Full rights notice: licenses/arxiv-2604.13700-CC-BY-4.0.txt.\par\smallskip

\noindent\textit{Editorial scope.} Counted unit: the Lemma whose source label is \texttt{lem:highoutandin} in the article (source0.tex lines 294--300, source label \texttt{lem:highoutandin}), delivered together with the article's own complete original proof of it (source0.tex lines 301--327, one \texttt{proof} environment of 27 lines). No mathematics is written, repaired or re-derived: statement and proof are the article's own text line for line. Required context retained uncounted: obs:densebound (lines 284--286). Every definition, display and label of those lines is retained unchanged. Omitted from this package and not redistributed: 1--283, 287--293, 328--424. Nothing inside a retained excerpt is omitted or altered: every retained body is delivered exactly as the article's lines are written, so no omission receipt of any kind accompanies this package. Citations: the retained excerpts carry no citation command at all, so no bibliography entry of the article is transcribed and no reference is redistributed. Reference adaptation: none was needed and none was made; every reference inside the delivered unit resolves inside it, so no unresolved reference or citation is produced. Printed slips kept literal: no printed word, symbol, hypothesis or number of the counted statement or of its proof was altered. Layout and class: the article's own class is \texttt{article} with packages including graphicx, geometry, amsthm, amssymb, amsmath, hyperref, xcolor, enumerate, listings, complexity, blkarray, lmodern, caption, subcaption; the wrapper is the lane's pinned \texttt{amsart} editorial wrapper at 10pt on A4, which restates the theorem-like environments the retained text needs and adds the article's own macro definitions that the retained text needs (none), with extra wrapper packages (bm, bbm, dsfont, xspace, cancel, mathtools). The delivered numbering is the unit's own. Assets: no retained span contains a figure environment, a graphics-inclusion command, a PDF-page inclusion, a table or a TikZ picture; no image file is copied into this package, the package declares no asset dependency and no image file is redistributed with it. Licence: version arXiv:2604.13700v2 of this article is distributed under the Creative Commons Attribution 4.0 International licence, as recorded in the arXiv licence metadata for this exact version and shown on the abstract page https://arxiv.org/abs/2604.13700v2. A full rights notice is delivered at licenses/arxiv-2604.13700-CC-BY-4.0.txt. Characters: every retained excerpt is pure ASCII, so the delivered engine, XeLaTeX with its default Latin Modern text font, reports no missing character.
\begin{observation}\label{obs:densebound}
For every digraph $D$, we have $\mathrm{a}(D)< \frac{\mathrm{v}(D)(\mathrm{v}(D)+\min\{c(D),2\mathrm{dtw}(D)\})}{2}$.
\end{observation}

\begin{lemma}\label{lem:highoutandin}
Let $r\in \mathbb{N}$, and let $\alpha,\beta,\gamma>0$ be such that $2\beta<\left(1-\frac{\alpha}{2}\right)^2$ and $$\gamma\ge \left(1-\frac{\alpha}{2}\right)-\sqrt{\left(1-\frac{\alpha}{2}\right)^2-2\beta}.$$ Then for every $(r,\beta,\gamma)$-dense digraph $D$ with $\Delta^+(D), \Delta^-(D)\le r$ and $\min\{c(D),2\mathrm{dtw}(D)\}<\alpha r$ we have $\mathrm{v}(D)\ge \left(\left(1-\frac{\alpha}{2}\right)+\sqrt{\left(1-\frac{\alpha}{2}\right)^2-2\beta}\right)r$.
Suppose further that $0<\delta<\min\{1/2,1-\sqrt{\beta}\}$ is such that
$$\left(1-\frac{\alpha}{2}\right)+\sqrt{\left(1-\frac{\alpha}{2}\right)^2-2\beta}\ge 2((1-\delta)-\sqrt{(1-\delta)^2-\beta}) $$ and
$$\frac{\beta}{1/2-\delta}\le 2((1-\delta)+\sqrt{(1-\delta)^2-\beta}).$$
Then $D$ contains a vertex $v$ such that $$\min\{d_D^+(v),d_D^-(v)\}\ge \delta r.$$
\end{lemma}

\begin{proof}
We start by proving the first part of the lemma concerning the number of vertices of $D$. Throughout the rest of this proof, it will be convenient to set $x:=\frac{\mathrm{v}(D)}{r}$, and our first goal is to show $x\ge \left(1-\frac{\alpha}{2}\right)+\sqrt{\left(1-\frac{\alpha}{2}\right)^2-2\beta}$. 

Recalling that $\min\{c(D),2\mathrm{dtw}(D)\}<\alpha r$ by assumption, Observation~\ref{obs:densebound} yields that
$$r\cdot \mathrm{v}(D)-\beta r^2\le \mathrm{a}(D)<\frac{\mathrm{v}(D)(\mathrm{v}(D)+\min\{c(D),2\mathrm{dtw}(D)\})}{2}<\frac{1}{2}\mathrm{v}(D)^2+\frac{\alpha r}{2}\mathrm{v}(D).$$

Dividing the previous inequality by $r^2$ we obtain:
$$x-\beta<\frac{1}{2}x^2+\frac{\alpha}{2}x.$$

Rearranging, we find that
$$x^2+(\alpha-2)x>-2\beta.$$
This can be further simplified to 
$$\left(x-\left(1-\frac{\alpha}{2}\right)\right)^2>\left(1-\frac{\alpha}{2}\right)^2-2\beta.$$ This implies that $x$ does not lie in the interval 
$$\left[\left(1-\frac{\alpha}{2}\right)-\sqrt{\left(1-\frac{\alpha}{2}\right)^2-2\beta},\left(1-\frac{\alpha}{2}\right)+\sqrt{\left(1-\frac{\alpha}{2}\right)^2-2\beta}\right].$$ Since $D$ is $(r,\beta,\gamma)$-dense, we have $x\ge \gamma\ge \left(1-\frac{\alpha}{2}\right)-\sqrt{\left(1-\frac{\alpha}{2}\right)^2-2\beta}$, and hence it follows that in fact, $x>\left(1-\frac{\alpha}{2}\right)+\sqrt{\left(1-\frac{\alpha}{2}\right)^2-2\beta}$, as desired. Hence $\mathrm{v}(D)>\left(\left(1-\frac{\alpha}{2}\right)-\sqrt{\left(1-\frac{\alpha}{2}\right)^2-2\beta}\right)r$, concluding the proof of the first part of the lemma.

Now suppose that $\delta\in (0,1/2)$ satisfies the inequalities stated in the second part of the lemma, and let us show that there exists a vertex $v$ in $D$ with out- and in-degree at least $\delta r$.
Towards a contradiction, suppose this is not the case, i.e. every vertex $v\in V(D)$ satisfies $d_D^+(v)<\delta r$ or $d_D^-(v)<\delta r$. Let $A, B$ be a partition of $V(D)$ such that $d_D^+(a)<\delta r$ for every $a\in A$ and $d_D^-(b)<\delta r$ for every $b\in B$. Since every arc in $D$ starts in $A$, or ends in $B$, or starts in $B$ and ends in $A$, we find that
$$r\cdot \mathrm{v}(D)-\beta r^2\le \mathrm{a}(D)\le \sum_{a\in A}d_D^+(a)+\sum_{b\in B}d_D^-(b)+\mathrm{a}_D(B,A)$$
$$<\delta r\cdot (|A|+|B|)+\min\{\Delta^+(D)\cdot |B|,\Delta^-(D)\cdot |A|\}$$
$$\le \delta r \cdot \mathrm{v}(D)+\frac{r\mathrm{v}(D)}{2}=\left(\delta+\frac{1}{2}\right)r\cdot \mathrm{v}(D).$$

Rearranging yields $\mathrm{v}(D)\le \frac{\beta}{\frac{1}{2}-\delta}r$, i.e. $x\le \frac{\beta}{\frac{1}{2}-\delta}$. In particular, $$x\in \left[\left(1-\frac{\alpha}{2}\right)+\sqrt{\left(1-\frac{\alpha}{2}\right)^2-2\beta},\frac{\beta}{\frac{1}{2}-\delta}\right].$$

Redoing the first steps of the previous estimate, but bounding $\mathrm{a}_D(B,A)$ now by the trivial upper bound $|A||B|\le \left(\frac{|A|+|B|}{2}\right)^2=\frac{1}{4}\mathrm{v}(D)^2$, we next obtain that
$$r\cdot \mathrm{v}(D)-\beta r^2\le \mathrm{a}(D)< \delta r\cdot \mathrm{v}(D)+\frac{1}{4}\mathrm{v}(D)^2.$$
Dividing by $r^2$, we obtain $x-\beta< \delta x+\frac{1}{4}x^2$. Rearranging now yields $(x-2(1-\delta))^2>4((1-\delta)^2-\beta)$. Taking roots, we obtain $|x-2(1-\delta)|>2\sqrt{(1-\delta)^2-\beta}$. This implies that $x$ does not lie in the interval $[2((1-\delta)-\sqrt{(1-\delta)^2-\beta}),2((1-\delta)+\sqrt{(1-\delta)^2-\beta})]$. This, however, is a contradiction, since we have previously shown that $$x\in \left[\left(1-\frac{\alpha}{2}\right)+\sqrt{\left(1-\frac{\alpha}{2}\right)^2-2\beta},\frac{\beta}{\frac{1}{2}-\delta}\right]\subseteq [2((1-\delta)-\sqrt{(1-\delta)^2-\beta}),2((1-\delta)+\sqrt{(1-\delta)^2-\beta})].$$ This shows that our above assumption was wrong: $D$ must indeed contain a vertex $v$ with out- and in-degree at least $\delta r$. This concludes the proof of the second part of the lemma. 
\end{proof}

\end{document}
